% Generated draft: requires mathematical and source review.
\documentclass[11pt]{article}
\usepackage[margin=25mm]{geometry}
\usepackage{fontspec}
\setmainfont{Latin Modern Roman}
\usepackage{amsmath,amssymb,mathtools,unicode-math}
\setmathfont{Latin Modern Math}
\usepackage{enumitem,needspace,xurl,hyperref}
\setlength{\parindent}{0pt}
\setlength{\parskip}{5pt}
\setlength{\emergencystretch}{3em}
\newcommand{\N}{\mathbb{N}}
\newcommand{\Z}{\mathbb{Z}}
\newcommand{\Q}{\mathbb{Q}}
\newcommand{\R}{\mathbb{R}}
\newcommand{\C}{\mathbb{C}}
\newcommand{\F}{\mathbb{F}}
\newcommand{\A}{\mathbb{A}}
\newcommand{\PP}{\mathbb{P}}
\newcommand{\E}{\mathbb{E}}
\newcommand{\eps}{\varepsilon}
\newcommand{\dd}{\,\mathrm{d}}
\begin{document}
\section*{UnsolvedMath ID: 40}
\textbf{Source problem code:} catalogue entry
\subsection{Definitions and mathematical statement}
For a spatial domain $\Omega$ in three dimensions, the incompressible inviscid Euler equations for velocity $u$ and pressure $p$ are
\[
\partial_t u+(u\cdot\nabla)u+\nabla p=0,\qquad \nabla\cdot u=0,\qquad u(0,x)=u_0(x).
\]
The source asks whether smooth initial velocity always gives a smooth solution for all forward times. It does not specify $\Omega$, the boundary conditions, or the smoothness and decay class, and these choices affect both the theorem and its status.

One precise conventional boundaryless reading is the following separate question for each $\Omega=\R^3$ and $\Omega=\mathbb T^3$, where $\mathbb T^3=\R^3/\Z^3$ is the periodic three-torus. Let $H^m(\Omega)$ denote the $L^2$ Sobolev space with weak derivatives through order $m$ in $L^2$. For every divergence-free
\[
u_0\in\bigcap_{m\geq0}H^m(\Omega;\R^3),
\]
does the corresponding classical Euler solution extend for every $t\geq0$, remaining smooth and belonging to $C([0,\infty);H^m(\Omega;\R^3))$ for every integer $m\geq0$? The velocity class has finite kinetic energy $\frac12\int_\Omega|u_0|^2\,dx$. This explicitly states a standard smooth Sobolev formulation; the source does not choose it over every other possible meaning of smooth initial data. A negative answer in that formulation would require such initial data whose classical solution loses smoothness in finite time, rather than merely a nonsmooth weak solution.

For a bounded domain, an initial-boundary-value problem must instead specify its domain and boundary condition, for example an impermeability condition $u\cdot n=0$ on a smooth boundary with outward unit normal $n$, together with suitable initial compatibility. That example is not silently imposed on the stored question.

The supplied literature triage dated 2026-08-17 reports that Chen and Hou construct a stable nearly self-similar finite-time singularity from smooth finite-energy axisymmetric three-dimensional Euler data in a smooth bounded cylinder with boundary, using analytic and rigorous numerical components. It therefore reports a negative answer to a universal formulation that includes the relevant bounded-domain model, while retaining the classical smooth-data questions on $\R^3$ and $\mathbb T^3$ as unresolved. The precise bounded-domain geometry and full theorem hypotheses are not supplied here and must be recovered from the cited papers before restating their theorem. The triage also warns that boundaryless results below the classical smooth threshold do not settle the smooth-data problem. Its partially-solved classification and the registry open label must be understood with this domain distinction.
\subsection{Short English statement}
Do smooth incompressible three-dimensional Euler flows stay smooth for all time? The domain must be specified: the supplied review reports a smooth-data singularity with boundary, while the standard whole-space and periodic smooth-data questions remain unresolved.
\subsection{Sources}
\begin{itemize}
\item Jiajie Chen and Thomas Y. Hou, ``Stable Nearly Self-Similar Blowup of the 2D Boussinesq and 3D Euler Equations with Smooth Data II: Rigorous Numerics,'' Multiscale Modeling \& Simulation 23 (2025): \url{https://doi.org/10.1137/23M1580395}.
\item Jiajie Chen and Thomas Y. Hou, ``Stable nearly self-similar blowup of the 2D Boussinesq and 3D Euler equations with smooth data I: Analysis,'' arXiv:2210.07191: \url{https://arxiv.org/abs/2210.07191}.
\item Jiajie Chen and Thomas Y. Hou, ``Singularity formation in 3D Euler equations with smooth initial data and boundary,'' Proceedings of the National Academy of Sciences 122 (2025), e2500940122: \url{https://doi.org/10.1073/pnas.2500940122}.
\item Local registry data/unsolvedmath/problems.json, permanent catalogue entry, exact source code catalogue entry, ``The Regularity Problem for Euler Equations''; statement, background, and literature triage dated 2026-08-17. This permanent record is distinct from ID 1515 despite the shared source code. Dataset attribution: Ulam AI and the UnsolvedMath Contributors.
\end{itemize}
\end{document}
